\documentclass[11pt]{article}
\usepackage{mathblatt}
% Präambel des Setzers (werkzeuge/setzer-praeambel.tex, wortgleich); dazu \gteil (Muster T2U)
\makeatletter
\setlength{\mbpfbe}{0mm}% keine Punkte (pruefung.md § 5: nur beim Prüfstein)
\newlength{\gEin}\setlength{\gEin}{0pt}
\let\gSchrift\relax
\renewcommand{\pfkreuzzeile}[1]{\par\noindent\hangindent3.2em\hangafter1\hspace*{1.6em}$\square$\ #1\par}
\newcommand{\gkreuz}[1]{$\square$\ #1\hspace{2em}}
\newcommand{\gkreuzzeile}[2]{\par\noindent\hangindent2.2em\hangafter1\mbox{#1\,$\square$}\ #2\par}% Buchstabe, Option: eine je Zeile (neben Skizze)
% links, was man ansieht, rechts, was man tut; Spaltengrenze überall an derselben Stelle (gemeinsam.md „Satz“)
% v0.17: Fußeintrag „Nr: Wert“ – kurz in einer Box (nie getrennt), lang umbrechbar mit der Nummer am Anfang
% (vorher \mbox: lange Lösungen liefen über den Rand, AS40 138 pt); robust, weil \fusshilfe in die .aux schreibt
\NewDocumentCommand{\gfuss}{m m}{\sbox0{#1:~#2}\ifdim\wd0<0.45\textwidth\mbox{#1:~#2}\else\mbox{#1:}\nobreak\ #2\fi}
\newsavebox{\gbox}\newlength{\gLinks}
% v0.10 (Lehrer 13.10.): keine starre Halbierung – die linke Spalte ist so breit wie die Figur (+ 6 mm), die
% Antwort steht direkt daneben; ist die Figur breiter als 0,62 Zeile, steht die Antwort darunter
\newcommand{\gzwei}[2]{\par\smallskip\noindent\sbox{\gbox}{#1}%
  \setlength{\gLinks}{\dimexpr\wd\gbox+6mm\relax}%
  \ifdim\gLinks>0.62\linewidth
    \usebox{\gbox}\par\smallskip\noindent\begin{minipage}{\linewidth}#2\end{minipage}%
  \else
    \begin{minipage}[t]{\gLinks}\vspace{0pt}\usebox{\gbox}\end{minipage}%
    \begin{minipage}[t]{\dimexpr\linewidth-\gLinks\relax}\vspace{0pt}\raggedright #2\end{minipage}%
  \fi\par}
% 16.10. (Lehrer, PA10 S. 1): Text links, Koordinatensystem rechts daneben; ist das System breiter als 0,58 Zeile,
% steht es wie bisher unter dem Text
% Passt das System knapp nicht mehr auf die Seite (bis 25 % zu hoch), wird es auf den Rest verkleinert statt
% die Seite halb leer zu lassen
\newlength{\gRest}\newlength{\gHoch}
\newcommand{\gtextks}[2]{\sbox{\gbox}{#1}%
  \setlength{\gHoch}{\dimexpr\ht\gbox+\dp\gbox\relax}%
  \ifdim\pagegoal<\maxdimen\setlength{\gRest}{\dimexpr\pagegoal-\pagetotal-4\baselineskip\relax}%
    \ifdim\gHoch>\gRest\ifdim\gRest>0.75\gHoch
      \sbox{\gbox}{\resizebox*{!}{\gRest}{\usebox{\gbox}}}\fi\fi\fi
  \ifdim\wd\gbox>0.58\linewidth #2\par\smallskip\noindent\gpassend{\usebox{\gbox}}\par
  \else\begin{minipage}[t]{\dimexpr\linewidth-\wd\gbox-5mm\relax}\raggedright #2\end{minipage}\hfill
    \begin{minipage}[t]{\wd\gbox}\vspace{0pt}\usebox{\gbox}\end{minipage}\par\fi}
% v0.9 (12.10., pruefung.md § 5 „Marke mit B, Bedingungen rechts“): drittes Argument von pfaufg = Bedingungen
% der Prüfung (v0.10: \pfohnetr = TR durchgestrichen, „FOR“), klein grau im rechten Rand auf der Grundlinie der Nummer, Zeilen mit \\
\newcommand{\gbedrand}[2]{\ifblank{#2}{}{\rlap{\hspace*{\dimexpr#1+2mm\relax}%
  \parbox[t]{15mm}{\raggedright\scriptsize\leavevmode\color{mbgrau}\strut#2}}}}
% Teil eines Bündels (a, b …): Marke links im Rand, Bedingungen rechts im Rand, auf der Zeile des Teils
\newcommand{\gteilrand}[2]{\llap{\mbpfrandmarke{#1}\hspace*{\mbpfnr}}\gbedrand{\linewidth}{#2}}
% Teil eines Bündels: Buchstabe in eigener Spalte, Inhalt eingerückt (\gEin = Einzug, damit die Spaltengrenze von
% \gzwei an derselben Stelle bleibt)
\newenvironment{gbteil}[3]{\par\smallskip\noindent\gteilrand{#2}{#3}\makebox[\mbpfteil][l]{\textbf{#1}}%
  \begin{minipage}[t]{\dimexpr\linewidth-\mbpfteil\relax}\setlength{\gEin}{\mbpfteil}\raggedright\strut\ignorespaces}%
 {\end{minipage}\par\addvspace{6pt}}% v0.18 (PG80 Nr. 1): Abstand, damit die Linie von a) nicht an b) rutscht
% Grafik höchstens so breit wie die Spalte (Bündel nebeneinander, Koordinatensystem ohne Karo)
\newsavebox{\gpbox}
\newcommand{\gpassend}[1]{\sbox{\gpbox}{#1}\ifdim\wd\gpbox>\linewidth\typeout{GPASSEND: verkleinert \the\wd\gpbox\space auf \the\linewidth}\resizebox{\linewidth}{!}{\usebox{\gpbox}}\else\usebox{\gpbox}\fi}
% Skizzenfeld (v0.10, Lehrer 13.10.): Karo ab dem linken Rand über die Zeile, „Skizze“ klein grau oben links im Feld
% (wie \pfskizzenfeld der Vorlage); der rechte Rand bleibt den Bedingungen
\newcommand{\gskizze}[1]{\par\smallskip\noindent\begin{tikzpicture}
  \clip (0,0) rectangle ({\linewidth-1mm},{#1*10mm});
  \draw[black!16,line width=0.3pt,step=5mm] (0,0) grid ({\linewidth},{#1*10mm});
  \node[anchor=north west,inner sep=2pt,font=\scriptsize,text=mbgrau,fill=white] at (0,{#1*10mm}) {Skizze};
  \end{tikzpicture}\par}
\RenewDocumentEnvironment{pfaufg}{m m m}{%
  \begin{lrbox}{\mbaufgabebox}\begin{minipage}[t]{\linewidth}%
  \gSchrift\noindent\hspace*{\gEin}\mbpfrandmarke{#1}\gbedrand{\linewidth-\gEin}{#3}%
  \begin{minipage}[t]{\mbpfnr}\strut\textbf{#2}\end{minipage}%
  \begin{minipage}[t]{\dimexpr\linewidth-\gEin-\mbpfnr\relax}\raggedright\strut\ignorespaces}%
 {\end{minipage}%
  \end{minipage}\end{lrbox}%
  \setlength{\mbaufgabehoehe}{\dimexpr\ht\mbaufgabebox+\dp\mbaufgabebox\relax}%
  \ifdefempty{\gkopfnach}{\par\addvspace{5pt}\Needspace*{\mbaufgabehoehe}}{%
    % Stufenkopf, der mit dieser (ersten) Aufgabe zusammenbleibt: Platz für beide verlangen, dann den Kopf setzen
    % (12.10., Quadratische: ein Kopf stand allein unten, die Aufgabe mit Koordinatensystem auf der nächsten Seite)
    % v0.17 (AS40 „Terme deuten“ allein unten): hinter dem Kopf steht das Sprungziel (\special, ein Knoten, nach
    % dem die Seite brechen durfte) – darum \nobreak nach dem Kopf; Vorhalt 56 pt = Kopf mit Linie und Abständen
    \par\Needspace*{\dimexpr\mbaufgabehoehe+56pt\relax}\gkopfnach\par\nobreak\gdef\gkopfnach{}\addvspace{5pt}}%
  \noindent\usebox{\mbaufgabebox}\par\addvspace{5pt}}
% v0.18 (plan.md § 10, 15.10. l): Abschnittsbuchstabe A, B, C vor dem Namen im Kopf – als Vorsatz, damit der
% Name im TeX gleich bleibt (Bestellkarte, Sprungziele lesen ihn)
\def\gbstakt{}\newcommand{\gbstvor}[1]{\gdef\gbstakt{#1}}
\NewCommandCopy\gpfstufekopfalt\pfstufekopf
\RenewDocumentCommand{\pfstufekopf}{o m m}{\IfValueTF{#1}{\gpfstufekopfalt[#1]}{\gpfstufekopfalt}%
  {\ifx\gbstakt\empty\else\gbstakt\enspace\fi #2}{#3}\gdef\gbstakt{}}
% v0.18 (16.10. e): Begriff grau, ohne Nummer, ohne Rahmen über der ersten Aufgabe der Stufe
\newcommand{\gbegriff}[1]{\par\nobreak\noindent\hspace*{\mbpfnr}\parbox[t]{\dimexpr\linewidth-\mbpfnr\relax}{\raggedright\small\color{mbgrau}#1}\par\nobreak\addvspace{3pt}}
% v0.18 (16.10. g): Lösungsblatt – Kopf mit kräftigem Strich, Abschnitt klein grau als Zwischenzeile
\newcommand{\glkopf}[1]{\noindent{\large\bfseries Lösungen \textperiodcentered{} #1}\par\vspace{3pt}\hrule height 1.2pt\par\medskip}
\newcommand{\glzwischen}[1]{\par\addvspace{8pt}\Needspace*{4\baselineskip}\noindent{\footnotesize\color{mbgrau}#1}\par\nobreak\vspace{1pt}}
% v0.18 (15.10. o, w): Blattende – Kern ohne Rahmen, dann graue Zeile „Mehr: n Aufgaben“
\newcommand{\gkern}[1]{\par\addvspace{14pt}\Needspace*{6\baselineskip}\noindent{\bfseries Kern}\par\nobreak\smallskip
  {\small\raggedright #1\par}}
\newcommand{\gmehr}[1]{\par\addvspace{10pt}\noindent{\footnotesize\color{mbgrau}Mehr: #1 Aufgaben}\par}
% Stufenkopf vormerken: \gkopfvor{<Kopf>} – die nächste pfaufg setzt ihn, sobald beide zusammen Platz haben
\def\gkopfnach{}
\newcommand{\gkopfvor}[1]{\gdef\gkopfnach{#1}}
% Teilaufgabe mit Skizze links in gzwei: Buchstabe in eigener Spalte, darunter Text und Skizze
\newcommand{\gteil}[2]{\makebox[\mbpfteil][l]{\textbf{#1}}\parbox[t]{\dimexpr0.5\textwidth-\gEin-\mbpfnr-5mm-\mbpfteil\relax}{\raggedright #2}}
% Stufenkopf der Vorlage (mathblatt.sty 2026-10-12: Linie, Abstand, 13 pt, \pfjahre) ohne ihren Vorhalt von
% 14 Zeilen: der Setzer fordert vor dem Kopf 0,16 Seitenhöhe (Kopf + erste Aufgabe), mehr nicht
\setlength{\mbpfstufevorhalt}{0pt}
\makeatother
\newcommand{\karo}[1]{\par\smallskip\noindent\begin{tikzpicture}
  \clip (0,0) rectangle ({\linewidth-1mm},{#1*10mm});
  \draw[black!16,line width=0.3pt,step=5mm] (0,0) grid ({\linewidth},{#1*10mm});
  \end{tikzpicture}\par}
\newcommand{\kk}[1]{$\square$\,#1\hspace{0.9em}}
\newcommand{\linie}[1]{\,{\color{black!45}\rule[-1pt]{#1}{0.4pt}}\,}
\newcommand{\pkt}{\textperiodcentered{}}
% v0.14 (14.10.): Stufenkopf mit fünf Punkten (Option = Jahre für die Bestellkarte) oder „zuletzt 19“
\NewDocumentCommand{\gpunkte}{o m}{#2}
\newcommand{\gzuletzt}[1]{zuletzt #1}
% v0.7 (12.10.): Schülerfrage unter dem Stufenkopf; Ankreuzoption als Bild (Buchstabe und Kästchen darunter)
\newcommand{\gfrage}[1]{\par\nobreak\noindent\hspace*{\mbpfnr}{\small\itshape #1}\par\nobreak\addvspace{2pt}}
\newlength{\gbildbreite}\setlength{\gbildbreite}{0.32\linewidth}% Setzer setzt sie je Reihe (vier Bilder: 0,24)
\newcommand{\gbildoption}[3]{\begin{minipage}[t][#3][b]{\gbildbreite}\centering #2\par\smallskip\mbox{#1\,$\square$}\end{minipage}}
% Lineare (12.10.): Bild mit Buchstabe ohne Kästchen (Diagramm zuordnen); Aussage mit „wahr / falsch“ zum Ankreuzen
\newcommand{\gbildname}[3]{\begin{minipage}[t][#3][b]{\gbildbreite}\centering #2\par\smallskip\mbox{\textbf{#1}}\end{minipage}}
\newcommand{\gwahrfalsch}[2]{\par\noindent\hangindent2.2em\hangafter1\makebox[2.2em][l]{(#1)}#2\hfill\mbox{wahr $\square$\quad falsch $\square$}\par}
\makeatletter\setlength{\mbpfbe}{0mm}\makeatother
\begin{document}

\pfheftstil
\pflaufkopf[P10 \pkt{} mehr]{Diagramme}{PS40-5}\pfseitevon
\pfheftkopf{Diagramme}{mehr}
\pfstufenzeile{Winkel gesucht \pkt{} Kreisdiagramm zeichnen \pkt{} Säule ergänzen \pkt{} Aussagen prüfen \pkt{} falschen Eindruck erklären}
% 2021-OS-K5b-m1 (Schwere I, rang 4)
\begin{pfaufg}{nach~P10\\’21~\pkt{}~5b}{1.}{}
Die Tabelle gibt an, welche Güter im Hamburger Hafen be- und entladen werden (Anteile in \%): Schrott 14, Flüssigkeiten ?, Container 62, Getreide 9, Sonstiges 2, gesamt 100. Ergänze den fehlenden Anteil für Flüssigkeiten. Schreib an jeden Sektor des Kreisdiagramms das passende Gut.
\par\smallskip\noindent\gpassend{\kreisdiagrammleer[1.8]{a/9,b/13,c/14,d/2,e/62}}\par
\karo{3}
\fusshilfe{\gfuss{1}{Flüssigkeiten 13 \%}}
\end{pfaufg}
% 2018-OS-K3c-m1 (Schwere II, rang 5)
\begin{pfaufg}{nach~P10\\’18~\pkt{}~3c}{2.}{}
Frauen wurden gefragt, wie oft sie Fleisch oder Wurst essen. Das Säulendiagramm zeigt die Anteile der Antworten. Stelle die Anteile aus dem Säulendiagramm als Streifendiagramm dar. Der vorgezeichnete Streifen ist 15 cm lang.
\par\smallskip\noindent\gpassend{\saeulenab[ymax=70,ystep=10,yfein=2,ylabel=Anteil in \%,breite=2.2]{nie/4,1- bis 3-mal pro Woche/12,4- bis 6-mal pro Woche/60,täglich/24}\par\medskip\noindent \begin{tikzpicture}[line width=0.6pt]\draw[black!20,line width=0.3pt,step=0.5cm] (-0.5,-0.5) grid (15.5,2);\draw[line width=0.8pt] (0,0) rectangle (15,1.5);\end{tikzpicture}}\par
\karo{2}
\fusshilfe{\gfuss{2}{0,6 cm nie, 1,8 cm 1–3-mal, 9,0 cm 4–6-mal, 3,6 cm täglich}}
\end{pfaufg}
% 2020-OS-K2a-m1 (Schwere I, rang 5)
\begin{pfaufg}{nach~P10\\’20~\pkt{}~2a}{3.}{}
Ein Schwimmbad zählt in der ersten Woche der Sommerferien jeden Tag seine Besucher. Am Donnerstag kamen 155 Besucher. Im Balkendiagramm fehlt der Balken für Donnerstag noch. Zeichne im Diagramm den Balken für Donnerstag ein.
\par\smallskip\noindent\gpassend{\balkenab[xmax=250,xstep=50,xfein=10,xlabel=Anzahl der Besucher]{So/240,Sa/210,Fr/170,Do/,Mi/180,Di/120,Mo/150}}\par
\karo{2}
\fusshilfe{\gfuss{3}{Balken Do bis 155}}
\end{pfaufg}
% 2017-OS-K2c-m1 (Schwere II, rang 6)
\begin{pfaufg}{nach~P10\\’17~\pkt{}~2c}{4.}{}
Die Tabelle zeigt, als wievieltes Kind ihrer Mutter die Kinder des Jahrgangs 2015 in Deutschland zur Welt kamen (Anteil in Prozent): erstes Kind 48,7 \%, zweites Kind 35,1 \%, drittes Kind 11,8 \%, viertes oder weiteres Kind 4,4 \%. Im Kreisdiagramm ist einer dieser Anteile grau gefärbt. Schreib an den grauen Sektor, zu welchem Anteil er gehört. Berechne, wie groß der Mittelpunktswinkel für „drittes Kind“ sein muss. Zeichne diesen Sektor in den Kreis ein und beschrifte ihn.
\par\smallskip\noindent\gpassend{\begin{tikzpicture}[line width=0.6pt]\filldraw[fill=black!30] (0,0) -- (180:2.0) arc (180:195.84:2.0) -- cycle;\draw (0,0) circle (2.0);\fill (0,0) circle (1.4pt);\draw[-{Stealth[length=2mm]},line width=0.5pt] (-3.3,-0.9) -- (187.92:1.50);\draw[line width=0.4pt] (-6.3,-0.90) -- (-3.3,-0.90);\draw[line width=0.4pt] (-6.3,-1.50) -- (-3.3,-1.50);\draw[line width=0.4pt] (-6.3,-2.10) -- (-3.3,-2.10);\end{tikzpicture}}\par
\karo{4}
\fusshilfe{\gfuss{4}{viertes oder weiteres Kind ($\approx$ 16°); 42°}}
\end{pfaufg}
% 2018-OS-K3a-m1 (Schwere II, rang 6)
\begin{pfaufg}{nach~P10\\’18~\pkt{}~3a}{5.}{}
1100 Jugendliche haben ihr Lieblingsessen genannt: Schnitzel 100, Pizza 130, Kartoffelgerichte 180, Gemüsegerichte 180, Fischgerichte 150; die Zahlen für Pasta und Salat fehlen noch. Die fünf bekannten Zahlen sind in einem Balkendiagramm dargestellt. Die Achse „Anzahl“ hat aber noch keine Einteilung, und die Balken sind nur mit A bis E bezeichnet. Trag an der Achse eine passende Einteilung ein. Schreib zu jedem Balken das passende Gericht.
\par\smallskip\noindent\gpassend{\begin{tikzpicture}[x=0.5cm,y=0.5cm,line width=0.6pt,font=\small]\draw[black!20,line width=0.3pt] (0,0) grid[step=1] (11,15);\filldraw[fill=black!15] (0,12) rectangle (7.5,14);\node[left] at (0,13) {A};\draw[line width=0.4pt] (12,12) -- (17,12);\filldraw[fill=black!15] (0,9) rectangle (5,11);\node[left] at (0,10) {B};\draw[line width=0.4pt] (12,9) -- (17,9);\filldraw[fill=black!15] (0,6) rectangle (6.5,8);\node[left] at (0,7) {C};\draw[line width=0.4pt] (12,6) -- (17,6);\filldraw[fill=black!15] (0,3) rectangle (9,5);\node[left] at (0,4) {D};\draw[line width=0.4pt] (12,3) -- (17,3);\filldraw[fill=black!15] (0,0) rectangle (9,2);\node[left] at (0,1) {E};\draw[line width=0.4pt] (12,0) -- (17,0);\foreach \x in {1,...,11}{\draw (\x,0.2) -- (\x,-0.2);}\draw[-{Stealth[length=2mm]}] (0,0) -- (11.8,0) node[below left,inner sep=3pt,yshift=-2pt] {Anzahl};\draw (0,0) -- (0,15);\end{tikzpicture}}\par
\karo{4}
\fusshilfe{\gfuss{5}{7,5 Kästchen}}
\end{pfaufg}
% 2018-OS-K3d-m1 (Schwere II, rang 7)
\begin{pfaufg}{nach~P10\\’18~\pkt{}~3d}{6.}{}
Frauen wurden gefragt, wie oft sie Fleisch oder Wurst essen. Das Säulendiagramm zeigt die Anteile der Antworten. Prüfe die beiden Aussagen zum Säulendiagramm. Kreuze jeweils an, ob sie richtig, falsch oder nicht entscheidbar ist, und begründe. Aussage 1: „Nicht einmal die Hälfte der Frauen isst 4- bis 6-mal pro Woche Fleisch oder Wurst.“ Aussage 2: „1- bis 3-mal pro Woche Fleisch oder Wurst isst jede 12. Frau.“
\par\smallskip\noindent\gpassend{\saeulenab[ymax=70,ystep=10,yfein=2,ylabel=Anteil in \%,breite=2.2]{nie/4,1- bis 3-mal pro Woche/12,4- bis 6-mal pro Woche/60,täglich/24}\par\medskip\noindent {\small\renewcommand{\arraystretch}{1.25}\begin{tabular}{|p{8.2cm}|c|c|c|}\hline Aussage & richtig & falsch & nicht entscheidbar\\\hline Aussage 1 &  &  & \\\hline Aussage 2 &  &  & \\\hline\end{tabular}}}\par
\karo{5}
\fusshilfe{\gfuss{6}{Aussage 1 falsch (60 \%); Aussage 2 falsch (12 \%)}}
\end{pfaufg}
% 2019-OS-K5c-m1 (Schwere III, rang 8)
\begin{pfaufg}{nach~P10\\’19~\pkt{}~5c}{7.}{}
2016 erschien ein Zeitungsartikel mit der Überschrift „Bücherlesen bald out?“. Darin heißt es: Die Buchhändler klagen, dass jedes Jahr dramatisch weniger Jugendliche Bücher lesen. Man fürchte, dass schon bald kein einziger Jugendlicher mehr zu einem Buch greift. Neben dem Text steht das Säulendiagramm. Schüler einer 10. Klasse halten zwei Aussagen des Artikels für falsch. Nenne eine davon und berichtige sie.
\par\smallskip\noindent\gpassend{\saeulenab[ymin=35,ymax=50,ystep=5,yfein=1,ylabel=Bücher lesende Jugendliche in \%]{2012/47,2013/45,2014/44,2015/42,2016/40}}\par
\karo{5}
\fusshilfe{\gfuss{7}{„extrem abnimmt“ falsch; „bald keiner mehr“ falsch}}
\end{pfaufg}
% 2014-OS-K3d-m1 (Schwere III, rang 9)
\begin{pfaufg}{nach~P10\\’14~\pkt{}~3d}{8.}{}
Eine andere Bank zahlt nur 0,4 \% Zinsen im Jahr. Für ihre Werbung „Tolle Zinsen!“ zeichnet sie ein Säulendiagramm, in dem das Guthaben kräftig zu wachsen scheint. Erkläre, mit welchem Trick die Grafik die Entwicklung des Guthabens besser aussehen lässt, als sie ist.
\par\smallskip\noindent\gpassend{\saeulenab[ymin=1040,ymax=1054,ystep=2,yfein=1,ylabel=Guthaben in Euro,breite=2.2]{nach einem Jahr/1044,nach zwei Jahren/1048,nach drei Jahren/1052}}\par
\karo{4}
\fusshilfe{\gfuss{8}{Zuwachs je Jahr nur 4 €}}
\end{pfaufg}
\par\vfill\noindent{\scriptsize\color{mbgrau}Originale: P10 2021 \pkt{} 5b, P10 2018 \pkt{} 3c, P10 2020 \pkt{} 2a, P10 2017 \pkt{} 2c, P10 2018 \pkt{} 3a, P10 2018 \pkt{} 3d, P10 2019 \pkt{} 5c, P10 2014 \pkt{} 3d\par}
\end{document}